\section{Lokalizacje, pochodne \texorpdfstring{$\infty$-kategorie}{kategorie wyższe} i teoria Mority}\label{cat:sec:derived-localization-morita}

W rozdziale~\ref{cat:sec:triangulated-derived} odwracaliśmy quasi-izomorfizmy w zwykłej kategorii, a w rozdziale~\ref{cat:sec:dg-a-infty} zachowaliśmy kompleksy morfizmów i opisaliśmy kategorię pochodną modułów DG. Teraz rozróżnimy lokalizację zwykłą od lokalizacji zachowującej przestrzenie odwzorowań i ich wyższe homotopie. Następnie będziemy badać pochodne kategorie modułów oraz równoważności Mority. Stabilność rozumiemy w sensie rozdziału~\ref{cat:sec:stable-infty-categories}. Ogólną teorię prezentowalnych $\infty$-kategorii i lokalizacji Bousfielda pozostawiamy do osobnego rozdziału.

\subsection{Kategorie względne i własność uniwersalna}

\begin{definicja}[Kategoria względna]\label{cat:def:relative-category}
Kategorią względną nazywamy parę $(\mathcal C,W)$, gdzie $\mathcal C$ jest kategorią, a $W$ jest jej szeroką podkategorią (zawiera wszystkie obiekty i identyczności), zawierającą wszystkie izomorfizmy. Morfizmy w $W$ nazywamy słabymi równoważnościami. Funktor względny $(\mathcal C,W)\to(\mathcal D,V)$ to funktor $F:\mathcal C\to\mathcal D$ taki, że $F(W)\subseteq V$.
\end{definicja}

Nie wymagamy w tej definicji warunku dwóch z trzech ani istnienia struktur faktoryzacji. Gdy $W$ jest klasą słabych równoważności kategorii modelowej, warunek dwóch z trzech zachodzi na mocy aksjomatów z definicji~\ref{cat:def:model-category}. Drugim przykładem jest kategoria kompleksów modułów z klasą quasi-izomorfizmów: składanie, identyczności i warunek dwóch z trzech wynikają z funktorialności kohomologii i odpowiednich własności izomorfizmów.

\begin{definicja}[Lokalizacja $\infty$-kategorii]\label{cat:def:infty-localization}
Niech $\mathcal C$ będzie $\infty$-kategorią, a $W$ zbiorem jej krawędzi. Przez $\operatorname{Fun}_W(\mathcal C,\mathcal E)$ oznaczamy pełną podkategorię $\infty$-kategorii funktorów, której obiektami są funktory przenoszące każdą krawędź z $W$ na równoważność. Lokalizacją $\mathcal C$ względem $W$ nazywamy funktor
\[
 L:\mathcal C\to\mathcal C[W^{-1}]_\infty
\]
taki, że dla każdej $\infty$-kategorii $\mathcal E$ prekompozycja z $L$ przyjmuje wartości w $\operatorname{Fun}_W(\mathcal C,\mathcal E)$ i indukuje równoważność
\begin{equation}\label{cat:eq:infty-localization-universal}
 L^*:\operatorname{Fun}(\mathcal C[W^{-1}]_\infty,\mathcal E)
 \xrightarrow{\ \simeq\ }\operatorname{Fun}_W(\mathcal C,\mathcal E).
\end{equation}
Dla kategorii względnej $(\mathcal C_0,W)$ przez jej lokalizację wyższą rozumiemy lokalizację nerwu $N\mathcal C_0$ względem odpowiadających krawędzi.
\end{definicja}

Wszystkie kwestie małości rozpatrujemy w ustalonych uniwersach; gdy odwracamy klasę morfizmów dużej kategorii, możemy najpierw powiększyć uniwersum. Pełność podkategorii funktorów oznacza, że nie narzucamy dodatkowych warunków na transformacje naturalne pomiędzy funktorami odwracającymi $W$. Własność~\eqref{cat:eq:infty-localization-universal} jest silniejsza od własności dotyczącej samych klas równoważności funktorów, ponieważ zawiera również informację o przestrzeniach transformacji.

\begin{twierdzenie}[Istnienie lokalizacji; wynik zewnętrzny]\label{cat:thm:infty-localization-exists}
Każda mała $\infty$-kategoria $\mathcal C$ i zbiór jej krawędzi $W$ mają lokalizację w sensie definicji~\ref{cat:def:infty-localization}. Lokalizację można wybrać jako quasi-kategorię z funktorem z $\mathcal C$.
\end{twierdzenie}

Jest to ogólne twierdzenie o lokalizacjach $\infty$-kategorii. Przyjmujemy je jako wynik zewnętrzny, nie jako konsekwencję konstrukcji zwykłej kategorii ilorazowej. Dokładne źródło: J. Lurie, \emph{Kerodon}, definicja 6.3.1.9 i stwierdzenie 6.3.2.1, \url{https://kerodon.net/tag/01ME}, \url{https://kerodon.net/tag/01MZ}. Poniżej wyprowadzamy konsekwencje z własności uniwersalnej, bez ponownego dowodzenia tego twierdzenia o istnieniu.

\begin{stwierdzenie}[Jednoznaczność lokalizacji]\label{cat:prop:infty-localization-unique}
Jeśli $L:\mathcal C\to\mathcal D$ i $L':\mathcal C\to\mathcal D'$ są lokalizacjami względem tej samej klasy $W$, to istnieje równoważność $F:\mathcal D\to\mathcal D'$ i naturalna równoważność $FL\simeq L'$. Przestrzeń wyborów pary $(F,FL\simeq L')$ jest ściągalna.
\end{stwierdzenie}
\begin{proof}
Z~\eqref{cat:eq:infty-localization-universal} dla $\mathcal E=\mathcal D'$ wynika, że $L'$ ma przeciwobraz przez równoważność $L^*$, czyli istnieje $F$ wraz z równoważnością $FL\simeq L'$. Zamieniając role lokalizacji, otrzymujemy $G:\mathcal D'\to\mathcal D$ z $GL'\simeq L$. Wobec tego
\[
 GFL\simeq GL'\simeq L,\qquad
 FGL'\simeq FL\simeq L'.
\]
Pełność i wierność w sensie homotopijnym funktorów prekompozycji pozwala podnieść te równoważności do $GF\simeq\id_{\mathcal D}$ i $FG\simeq\id_{\mathcal D'}$. Dokładniej, podnosimy krawędzie na mocy równoważności przestrzeni odwzorowań i otrzymujemy krawędzie odwracalne, ponieważ funktor pełny i wierny w sensie homotopijnym odbija równoważności: odpowiednią odwrotność i obie równości z identycznościami podnosimy na poziomie kategorii homotopii. Zatem $F$ jest równoważnością.

Przechodząc w~\eqref{cat:eq:infty-localization-universal} do rdzeni (definicja~\ref{cat:def:quasicat-core}), otrzymujemy równoważność przestrzeni. Homotopijne włókno równoważności przestrzeni nad punktem $L'$ jest ściągalne: równoważność identyfikuje je z włóknem identyczności, czyli z przestrzenią dróg kończących się w tym punkcie, która jest ściągalna. To włókno jest właśnie przestrzenią par $(F,FL\simeq L')$. Każdy jego punkt odpowiada równoważności przez pierwszą część dowodu.
\end{proof}

\subsection{Lokalizacja etapami i nasycenie}

\begin{stwierdzenie}[Lokalizacja etapami]\label{cat:prop:infty-localization-iterated}
Jeżeli $W,V$ są zbiorami krawędzi $\mathcal C$, to lokalizacja $\mathcal C$ względem $W\cup V$ jest równoważna lokalizacji $\mathcal C[W^{-1}]_\infty$ względem obrazów krawędzi z $V$.
\end{stwierdzenie}
\begin{proof}
Niech $L_W:\mathcal C\to\mathcal D$ będzie pierwszą lokalizacją, a $L_V:\mathcal D\to\mathcal P$ lokalizacją względem $L_W(V)$. Dla każdej $\mathcal E$ prekompozycja z $L_V$ identyfikuje $\operatorname{Fun}(\mathcal P,\mathcal E)$ z pełną podkategorią $\operatorname{Fun}(\mathcal D,\mathcal E)$ na funktorach odwracających $L_W(V)$. Równoważność $L_W^*$ identyfikuje $\operatorname{Fun}(\mathcal D,\mathcal E)$ z $\operatorname{Fun}_W(\mathcal C,\mathcal E)$. Pod jej działaniem warunek odwracania $L_W(V)$ staje się dokładnie warunkiem odwracania $V$. Warunek ten jest niezmienniczy względem równoważności funktorów, gdyż naturalna równoważność daje na każdej krawędzi kwadrat przemienny w kategorii homotopii z izomorficznymi pionowymi strzałkami. Stąd ograniczenie drugiej równoważności jest równoważnością odpowiednich pełnych podkategorii. Ich złożenie identyfikuje $\operatorname{Fun}(\mathcal P,\mathcal E)$ z $\operatorname{Fun}_{W\cup V}(\mathcal C,\mathcal E)$ i jest prekompozycją z $L_VL_W$. Stosujemy stwierdzenie~\ref{cat:prop:infty-localization-unique}.
\end{proof}

\begin{definicja}[Nasycenie względem lokalizacji]\label{cat:def:infty-localization-saturation}
Dla lokalizacji $L$ klasy $W$ oznaczamy
\[
 \overline W=\{f\in\mathcal C_1:L(f)\text{ jest równoważnością}\}.
\]
Nazywamy $\overline W$ nasyceniem klasy $W$ względem lokalizacji. Nie utożsamiamy go z innymi pojęciami nasycenia, w szczególności nie wymagamy tu domknięcia na wszystkie kobazy czy kogranice.
\end{definicja}

\begin{stwierdzenie}[Własności nasycenia]\label{cat:prop:infty-localization-saturation}
Zachodzi $W\subseteq\overline W$. Nasycenie zawiera równoważności, jest niezmiennicze względem homotopii krawędzi, spełnia warunek dwóch z trzech oraz jest domknięte na retrakty w kategorii strzałek. Lokalizacja względem $\overline W$ jest równoważna lokalizacji względem $W$. Ponadto $\overline W$ jest dokładnie klasą krawędzi odwracanych przez każdy funktor odwracający $W$.
\end{stwierdzenie}
\begin{proof}
Pierwsza inkluzja wynika z definicji lokalizacji, biorąc funktor identyczności jej przeciwdziedziny. Funktor zachowuje równoważności, a krawędź jest równoważnością dokładnie wtedy, gdy jej klasa w kategorii homotopii jest izomorfizmem (definicja~\ref{cat:def:quasicat-equivalence}). Stąd niezmienniczość względem homotopii oraz warunek dwóch z trzech wynikają z odpowiednich własności izomorfizmów.

Dla retraktów wystarczy następujący rachunek w zwykłej kategorii: niech $f:X\to Y$ będzie retraktem izomorfizmu $g:U\to V$ w kategorii strzałek, z morfizmami włożenia $i_X,i_Y$ i retrakcji $r_X,r_Y$. Mamy $g i_X=i_Yf$, $fr_X=r_Yg$ oraz $r_Xi_X=\id_X$, $r_Yi_Y=\id_Y$. Wówczas $r_X g^{-1}i_Y$ jest odwrotnością $f$, gdyż
\[
 r_X g^{-1}i_Y f=r_X i_X=\id_X,\qquad
 f r_Xg^{-1}i_Y=r_Yi_Y=\id_Y.
\]
Stosując ten rachunek w kategorii homotopii przeciwdziedziny $L$ otrzymujemy domknięcie na retrakty; funktor indukowany przez $L$ przenosi kwadraty i równości retrakcji na odpowiednie kwadraty i równości klas homotopii.

Jeśli $F:\mathcal C\to\mathcal E$ odwraca $W$, to istnieje $\widetilde F$ z $\widetilde FL\simeq F$ przez~\eqref{cat:eq:infty-localization-universal}. Zatem $F$ odwraca również $\overline W$, bo $\widetilde F$ zachowuje równoważności, a naturalna równoważność funktorów zachowuje warunek odwracania krawędzi. Odwrotna inkluzja jest oczywista przez $W\subseteq\overline W$. W szczególności pełne podkategorie funktorów odwracających $W$ i $\overline W$ są równe. To pokazuje, że $L$ spełnia również własność uniwersalną dla $\overline W$. Ostatnia teza wynika z powyższej faktoryzacji oraz z wyboru samego $L$ jako jednego z funktorów odwracających $W$.
\end{proof}

\subsection{Porównanie z lokalizacją zwykłą}

\begin{lemat}[Funktory do nerwu zwykłej kategorii]\label{cat:lem:infty-localization-ordinary-target}
Dla quasi-kategorii $\mathcal C$ i zwykłej kategorii $\mathcal E$ zachodzi kanoniczny izomorfizm zbiorów symplicjalnych
\[
 \operatorname{Fun}(\mathcal C,N\mathcal E)\cong
 N\operatorname{Fun}(h\mathcal C,\mathcal E),
\]
gdzie $h\mathcal C$ oznacza kategorię homotopii.
\end{lemat}
\begin{proof}
Odwzorowanie symplicjalne $\mathcal C\to N\mathcal E$ określa funktor $h\mathcal C\to\mathcal E$. Istotnie, obrazy krawędzi homotopijnych są równe, gdyż obrazy trójkątów wyznaczających homotopię są trójkątami nerwu z krawędzią będącą identycznością. Obraz trójkąta wyznaczającego złożenie pokazuje zachowanie składania. Odwrotnie, dla każdego $n$ sympleks $\sigma:\Delta^n\to\mathcal C$ wyznacza funktor $[n]\to h\mathcal C$ przez swoje wierzchołki i klasy krawędzi: relacje składania wynikają z jego ścian dwuwymiarowych. Złożenie z danym funktorem $h\mathcal C\to\mathcal E$ wyznacza $n$-sympleks $N\mathcal E$. Konstrukcja komutuje z operatorami symplicjalnymi. Obie konstrukcje są wzajemnie odwrotne, ponieważ sympleks nerwu jest wyznaczony przez ciąg kolejnych krawędzi.

W stopniu $n$ lewa strona lematu jest zbiorem map $\mathcal C\times\Delta^n\to N\mathcal E$. Produkt quasi-kategorii ma kategorię homotopii równą produktowi kategorii homotopii: wierzchołki i krawędzie są parami, a homotopie oraz wypełnienia trójkątów odpowiadających składaniu są wybierane składowo. Ponadto $h\Delta^n=[n]$ przez wniosek~\ref{cat:cor:homotopy-nerve}. Z poprzedniej części dowodu dostajemy bijekcję
\[
 \operatorname{Hom}_{\mathrm{sSet}}(\mathcal C\times\Delta^n,N\mathcal E)
 \cong\operatorname{Fun}(h\mathcal C\times[n],\mathcal E)_0
 \cong\operatorname{Fun}([n],\operatorname{Fun}(h\mathcal C,\mathcal E))_0.
\]
Indeks $0$ oznacza zbiór obiektów kategorii funktorów. Ostatnia bijekcja to zwykłe sprzężenie iloczynu i kategorii funktorów: funktor na produkcie daje rodzinę funktorów oraz transformacje naturalne wzdłuż strzałek $[n]$. Bijekcje są zgodne z operatorami symplicjalnymi, gdyż pochodzą z naturalnej prekompozycji w zmiennej $[n]$. Otrzymujemy żądany izomorfizm.
\end{proof}

\begin{twierdzenie}[Kategoria homotopii lokalizacji]\label{cat:thm:infty-localization-homotopy}
Dla lokalizacji $L:\mathcal C\to\mathcal D=\mathcal C[W^{-1}]_\infty$ funktor
\[
 hL:h\mathcal C\to h\mathcal D
\]
realizuje zwykłą lokalizację kategorii $h\mathcal C$ względem klas $[W]$. W szczególności dla zwykłej kategorii względnej zachodzi równoważność
\[
 h((N\mathcal C_0)[W^{-1}]_\infty)\simeq \mathcal C_0[W^{-1}].
\]
\end{twierdzenie}
\begin{proof}
W~\eqref{cat:eq:infty-localization-universal} bierzemy $\mathcal E=N\mathcal B$ dla dowolnej zwykłej kategorii $\mathcal B$. Z lematu~\ref{cat:lem:infty-localization-ordinary-target} otrzymujemy równoważność nerwów
\[
 N\operatorname{Fun}(h\mathcal D,\mathcal B)
 \simeq N\operatorname{Fun}_{[W]}(h\mathcal C,\mathcal B).
\]
Warunek odwracania krawędzi odpowiada tu odwracaniu klas, ponieważ równoważności w nerwie zwykłej kategorii są dokładnie jej izomorfizmami. Przechodząc do kategorii homotopii, dostajemy równoważność zwykłych kategorii funktorów indukowaną przez $hL$. Jest to własność uniwersalna zwykłej lokalizacji w wersji nieścisłej (z dokładnością do naturalnego izomorfizmu). Porównując ją z lokalizacją przez formalne odwrotności, otrzymujemy równoważne kategorie: podnosimy oba funktory lokalizacyjne i sprawdzamy, że złożenia są naturalnie izomorficzne z identycznościami, jak w dowodzie stwierdzenia~\ref{cat:prop:infty-localization-unique}. Dla zwykłej $\mathcal C_0$ używamy jeszcze $hN\mathcal C_0\cong\mathcal C_0$.
\end{proof}

\begin{przyklad}[Przestrzeń homotopii nie jest kategorią homotopii]\label{cat:ex:infty-localization-higher-information}
Niech $K$ będzie kompleksem Kana modelującym sferę $S^2$, na przykład zbiorem symplicjalnym singularnym sfery. Każda jego krawędź jest równoważnością, więc identyczność $K\to K$ realizuje lokalizację względem wszystkich krawędzi. Kategoria $hK$ jest równoważna kategorii końcowej, gdyż $S^2$ jest spójna drogowo i jednospójna. Mimo to $K$ nie jest równoważny punktowi jako $\infty$-kategoria: jako kompleks Kana ma nietrywialną grupę $\pi_2(K)\cong\mathbb Z$. Równoważność $\infty$-kategorii indukuje równoważność ich rdzeni, a rdzeń kompleksu Kana jest nim samym. Zwykła kategoria homotopii nie rozróżnia zatem punktu i $K$, natomiast lokalizacja wyższa je rozróżnia.
\end{przyklad}

\subsection{Konstrukcja przez odwrócenie pojedynczych strzałek}

Niech $\mathcal I$ będzie grupoidem o dwóch obiektach $0,1$ i dokładnie jednym morfizmie dla każdej uporządkowanej pary obiektów. Oznaczamy $J=N\mathcal I$. Włożenie $\Delta^1\to J$ wysyła krawędź $0\to1$ na jedyny izomorfizm o tych końcach.

\begin{lemat}[Odwrócenie strzałki w dowolnej $\infty$-kategorii; wynik zewnętrzny]\label{cat:lem:infty-localization-walking-iso}
Prekompozycja z $\Delta^1\to J$ indukuje równoważność
\[
 \operatorname{Fun}(J,\mathcal E)\simeq
 \operatorname{Fun}_{\mathrm{eq}}(\Delta^1,\mathcal E),
\]
gdzie po prawej bierzemy pełną podkategorię na krawędziach będących równoważnościami w $\mathcal E$.
\end{lemat}

Jest to szczególny przypadek twierdzenia o lokalizacji względem wszystkich krawędzi, przyjmowany tutaj jako wynik zewnętrzny. Zob. J. Lurie, \emph{Kerodon}, podrozdział 6.3.1, w szczególności stwierdzenie 6.3.1.21 (\url{https://kerodon.net/tag/01MY}). Rzeczywiście $J$ jest kompleksem Kana, a $\Delta^1\to J$ jest słabą równoważnością homotopijną, ponieważ oba zbiory symplicjalne są ściągalne: dla $J$ wynika to z faktu, że $\mathcal I$ ma obiekt końcowy (i początkowy).

\begin{stwierdzenie}[Opis lokalizacji jako kokartezjańskiego kwadratu]\label{cat:prop:infty-localization-pushout}
W $\infty$-kategorii małych $\infty$-kategorii $\mathrm{Cat}_\infty$ niech dany będzie kwadrat kokartezjański
\[
\begin{tikzpicture}[>=Stealth,baseline=(a.base)]
 \node (a) at (0,1.8) {$\displaystyle\coprod_{w\in W}\Delta^1$};
 \node (b) at (5,1.8) {$\displaystyle\coprod_{w\in W}J$};
 \node (c) at (0,0) {$\mathcal C$};
 \node (d) at (5,0) {$\mathcal P$};
 \draw[->] (a) -- (b);
 \draw[->] (a) -- node[left] {$w$} (c);
 \draw[->] (b) -- (d);
 \draw[->] (c) -- node[below] {$L$} (d);
\end{tikzpicture}
\]
w którym lewy morfizm jest zadany przez krawędzie $w$. Wtedy $L$ jest lokalizacją względem $W$.
\end{stwierdzenie}
\begin{proof}
Korzystamy z własności uniwersalnej kwadratu kokartezjańskiego w $\mathrm{Cat}_\infty$ i z lematu~\ref{cat:lem:infty-localization-walking-iso}. Najpierw na poziomie przestrzeni odwzorowań otrzymujemy homotopijny iloczyn włóknisty
\[
 \operatorname{Fun}(\mathcal P,\mathcal E)^\simeq
 \simeq
 \operatorname{Fun}(\mathcal C,\mathcal E)^\simeq
 \times_{\prod_w\operatorname{Fun}(\Delta^1,\mathcal E)^\simeq}
 \prod_w\operatorname{Fun}(J,\mathcal E)^\simeq.
\]
Każda składowa prawej mapy identyfikuje się z włączeniem sumy składowych spójności odpowiadających równoważnościom. Rzeczywiście warunek bycia równoważnością jest niezmienniczy względem naturalnej równoważności diagramów strzałek. Zatem iloczyn włóknisty identyfikuje się z $\operatorname{Fun}_W(\mathcal C,\mathcal E)^\simeq$.

Aby uzyskać własność dla całych kategorii funktorów, a nie tylko ich rdzeni, stosujemy ten sam argument do $\operatorname{Fun}(\mathcal K,\mathcal E)$ dla każdej małej $\mathcal K$. Używamy tu domkniętości kartezjańskiej $\mathrm{Cat}_\infty$ (standardowy wynik zewnętrzny o kategoriach funktorów): równoważnie $\operatorname{Fun}(\mathcal P,\operatorname{Fun}(\mathcal K,\mathcal E))\simeq\operatorname{Fun}(\mathcal K,\operatorname{Fun}(\mathcal P,\mathcal E))$. Równoważności w kategorii funktorów są punktowe (twierdzenie~\ref{cat:thm:quasicat-pointwise-equivalence}). Dlatego warunek odwracania $W$ po zamianie zmiennych oznacza dokładnie faktoryzację funktora $\mathcal K\to\operatorname{Fun}(\mathcal C,\mathcal E)$ przez pełną podkategorię $\operatorname{Fun}_W$. Otrzymujemy dla każdego $\mathcal K$ równoważność przestrzeni odwzorowań z $\mathcal K$ indukowaną przez
\[
 L^*:\operatorname{Fun}(\mathcal P,\mathcal E)
 \longrightarrow\operatorname{Fun}_W(\mathcal C,\mathcal E).
\]
Taki funktor jest równoważnością: biorąc jako $\mathcal K$ jego przeciwdziedzinę, podnosimy identyczność do odwrotności z prawej strony; biorąc jako $\mathcal K$ jego dziedzinę, z injektywności na składowych spójności dostajemy również odwrotność z lewej strony. W ten sposób spełniona jest~\eqref{cat:eq:infty-localization-universal}.
\end{proof}

\begin{uwaga}[Zakres konstrukcji]\label{cat:rem:infty-localization-pushout}
Powyższy kwadrat jest kwadratem kokartezjańskim w $\mathrm{Cat}_\infty$, a nie dowolnym ścisłym koproduktem amalgamowanym zbiorów symplicjalnych. Istnienie wymaganych kogranic w $\mathrm{Cat}_\infty$ oraz domkniętość kartezjańska są ogólnymi wynikami teorii $\infty$-kategorii; zob. J. Lurie, \emph{Kerodon}, podrozdział 6.3.2 (\url{https://kerodon.net/tag/01MZ}), zwłaszcza stwierdzenie 6.3.2.4 dotyczące kwadratów kokartezjańskich i lokalizacji. W modelu symplicjalnym należy zadbać o obliczenie właściwego homotopijnego kwadratu kokartezjańskiego i o zastąpienie wyniku quasi-kategorią. Nie wnioskujemy, że każdy morfizm lokalizacji wyższej jest jednoznacznie zadany pojedynczym ułamkiem.
\end{uwaga}

